% ===================================================================
%  CUADRO N.º 2 — El Cuerpo Humano. — El Recreo. — Los Juegos.
%  Traduction 2026 d'apres texto/io, controlee sur texto/fr et
%  texto/en. Consigne : voir texto/es/10-cuadro-01.tex.
% ===================================================================

%%K t02-apar-1 apar
\VUsaut{4.0mm}
\VUtitre{15.0pt}{60}{CUADRO N.º 2}
\VUsaut{2.0mm}
\VUtitre{11.4pt}{0}{El Cuerpo Humano. --- El Recreo.}
\VUtitre{11.4pt}{0}{Los Juegos.}

%%K t02-tit-0 sub
\VUtitre{13.2pt}{0}{El Cuerpo Humano.}
\VUsaut{2.4mm}
\VUcentre{10.2pt}{0}{\textit{(Lección sencilla de ciencias naturales.)}}
\VUsaut{3.0mm}

%%K t02-01-1 p
1. --- Las principales partes del cuerpo humano son: la \VUgras{cabeza}, el
\VUgras{tronco} y los \VUgras{miembros}.

%%K t02-tit-1 sub
\VUpk{10.2pt}{40}{La Cabeza.}
\VUsaut{3.0mm}

%%K t02-02-1 p
2. --- La cabeza consta de dos partes: el \VUgras{cráneo}, que contiene el
\VUgras{cerebro} y que cubre el \VUgras{pelo}, y la \VUgras{cara}.

%%K t02-03-1 p
3. --- En nuestro dibujo no vemos ni el cráneo ni el cerebro; pero vemos
con toda claridad las partes importantes de la cara.

%%K t02-04-1 p
4. --- He aquí primero la \VUgras{frente}, que anuncia una inteligencia
poderosa, un pensamiento siempre despierto. Las \VUgras{sienes} están muy
despejadas. El \VUgras{ojo} es vivo y bien abierto. Una \VUgras{ceja} espesa
y unas \VUgras{pestañas} largas lo protegen.

%%K t02-05-1 p
5. --- He aquí además la \VUgras{nariz} y las \VUgras{ventanas de la nariz}.
La nariz nos sirve para oler y para respirar. A cada lado de la nariz
están las \VUgras{mejillas}, y más allá las \VUgras{orejas}. La parte
inferior de la cara la forman la \VUgras{boca} y la \VUgras{barbilla}.

%%K t02-06-1 p
6. --- No las vemos muy bien, porque el \VUgras{bigote} y la \VUgras{barba}
nos las ocultan. Pero sabemos que la boca tiene los dos \VUgras{labios}, la
\VUgras{mandíbula superior} y la \VUgras{mandíbula inferior} con sus
\VUgras{dientes}, y la \VUgras{lengua}. Con la boca hablamos y comemos. Con
la lengua saboreamos los alimentos.

%%K t02-07-1 p
7. --- El ojo, la oreja, la nariz y la boca son, junto con los dedos, los
órganos de los cinco sentidos: la \VUgras{vista}, el \VUgras{oído}, el
\VUgras{olfato}, el \VUgras{gusto} y el \VUgras{tacto}.

%%K t02-08-1 p
8. --- La cabeza es, pues, una de las partes más interesantes del cuerpo
humano.

%%K t02-tit-2 sub
\VUsaut{4.4mm}
\VUpk{10.2pt}{40}{El Tronco.}
\VUsaut{3.0mm}

%%K t02-09-1 p
9. --- El \VUgras{cuello} une la cabeza al tronco. Comprende la
\VUgras{garganta} y la \VUgras{nuca}.

%%K t02-10-1 p
10. --- El tronco contiene también muchos órganos esenciales. En el
\VUgras{pecho} están los \VUgras{pulmones}, con los que respiramos, y el
\VUgras{corazón}, que es el centro de la vida. Cuando el corazón deja de
latir, el ser vivo muere.

%%K t02-11-1 p
11. --- Las \VUgras{costillas} sostienen el pecho.

%%K t02-12-1 p
12. --- Las demás partes del tronco son el \VUgras{costado}, la
\VUgras{espalda}, los \VUgras{hombros} y el \VUgras{abdomen} (vientre). Nos
acostamos de lado o de espaldas para dormir, y llevamos las cargas sobre
el hombro.

%%K t02-13-1 p
13. --- En el abdomen están el \VUgras{estómago} y los \VUgras{intestinos}.
Nos sirven para digerir los alimentos.

%%K t02-tit-3 sub
\VUsaut{4.4mm}
\VUpk{10.2pt}{40}{Los Miembros.}
\VUsaut{3.0mm}

%%K t02-14-1 p
14. --- Tenemos cuatro \VUgras{miembros}: dos \VUgras{brazos} y dos
\VUgras{piernas}. Con los brazos cogemos, sostenemos, levantamos y
recogemos los objetos; con las piernas nos tenemos de pie, andamos y
corremos.

%%K t02-15-1 p
15. --- El \VUgras{hombro} une el brazo al cuerpo. Un \VUgras{músculo} muy
robusto, el bíceps, mueve el brazo, que el \VUgras{codo} une al
\VUgras{antebrazo}. La \VUgras{muñeca} forma la articulación de la
\VUgras{mano}, que termina en los \VUgras{dedos} y las \VUgras{uñas}. En la
mano distinguimos una \VUgras{vena} (y una arteria) por la que corre la
\VUgras{sangre}.

%%K t02-16-1 p
16. --- Las partes de la pierna son: el \VUgras{muslo}, la \VUgras{rodilla} y
la \VUgras{corva}, la \VUgras{pantorrilla}, cuya \VUgras{carne} cubre la
\VUgras{piel}, el \VUgras{tobillo} y el \VUgras{pie}. Los \VUgras{huesos} de la
pierna son muy gruesos y muy fuertes. En el extremo del pie están los
\VUgras{dedos del pie}. Las demás partes son la \VUgras{planta} y el
\VUgras{talón}.

%%K t02-17-1 p
17. --- Tal es la descripción casi completa del cuerpo humano.

%%K t02-17-2 p
\textit{Nota.} --- \textit{Con ayuda de la expresión «me duele... (la
cabeza, etc.)», el profesor podrá repasar todo este vocabulario desde el
punto de vista del buen o del mal} \VUgras{estado del cuerpo}.

%%K t02-tit-4 sub
\VUsaut{5.0mm}
\VUpk{10.2pt}{40}{Gimnasia.}
\VUsaut{2.4mm}
\VUcentre{10.2pt}{0}{\textit{(Contado por uno de los alumnos.)}}
\VUsaut{3.0mm}

%%K t02-18-1 p
18. --- Para ser flexibles, ágiles y sanos, tenemos que ejercitar las
piernas y los brazos. Por eso jugamos y hacemos \VUgras{gimnasia}.

%%K t02-19-1 p
19. --- Entre semana, estos dos ejercicios se hacen en el patio del
instituto (véase el cuadro n.º 1). Pero los jueves y los domingos vamos a
una gran pradera, donde tenemos sitio para correr y divertirnos.

%%K t02-20-1 p
20. --- Nuestros maestros y nuestros padres nos acompañan allí a veces;
pero es el \VUgras{profesor de gimnasia} \textsuperscript{(12)} quien dirige los
ejercicios.

%%K t02-21-1 p
21. --- En la pradera, a la izquierda de la entrada, están los
\VUgras{aparatos de gimnasia} \textsuperscript{(1)}: trepamos por la
\VUgras{escalera} \textsuperscript{(2)}, hacemos el pino en las \VUgras{anillas}
\textsuperscript{(3)} y en el \VUgras{trapecio} \textsuperscript{(4)}, nos izamos con las
manos hasta lo alto de la \VUgras{escala de cuerda} \textsuperscript{(5)} o de la
\VUgras{cuerda de nudos} \textsuperscript{(6)}.

%%K t02-22-1 p
22. --- Mientras tanto, nuestros compañeros saltan por encima del
\VUgras{potro de madera} \textsuperscript{(7)} o de las \VUgras{barras paralelas}
\textsuperscript{(11)}. Otros \VUgras{boxean} \textsuperscript{(8)} o hacen
\VUgras{esgrima} \textsuperscript{(9)} con careta y \VUgras{florete}. Otros, en fin,
levantan \VUgras{pesas} \textsuperscript{(10)}.

%%K t02-tit-5 sub
\VUsaut{4.6mm}
\VUpk{10.2pt}{40}{Los Juegos.}
\VUsaut{3.0mm}

%%K t02-23-1 p
23. --- Alrededor de los gimnastas, niños y niñas juegan a distintos
\VUgras{juegos}. Las niñas se divierten con su \VUgras{aro} \textsuperscript{(14)};
saltan a la \VUgras{comba} \textsuperscript{(13)}, o invitan a sus hermanos y a sus
primos a una partida de \VUgras{croquet} \textsuperscript{(15)}. ¡Les encanta
apartar la \VUgras{bola} del contrario con su \VUgras{mazo de madera},
justo cuando este creía pasar sin trabajo por el arquillo!

%%K t02-24-1 p
24. --- Los niños prefieren los juegos más recios. Juegan de buena gana a
los \VUgras{bolos} \textsuperscript{(16)}, que derriban con maña con la
\VUgras{bola} \textsuperscript{(17)}. Tampoco desdeñan la \VUgras{peonza}
\textsuperscript{(18)} y el \VUgras{boliche} \textsuperscript{(19)}, ni siquiera el
\VUgras{juego de la rana} \textsuperscript{(20)}. Pero sus juegos preferidos son
las \VUgras{canicas} \textsuperscript{(21)} y la \VUgras{pelota a mano}
\textsuperscript{(22)}, porque la pared del \VUgras{invernadero} \textsuperscript{(26)} va
muy bien para hacer rebotar la pelota ligera.

%%K t02-25-1 p
25. --- El pequeño Juan, subido en sus \VUgras{zancos} \textsuperscript{(24)}, es
muy diestro y sabe guardar muy bien el equilibrio. Cerca de él, algunos
compañeros juegan al \VUgras{escondite} \textsuperscript{(25)} en ese invernadero y
detrás del \VUgras{seto}. Otros prefieren la \VUgras{gallina ciega con
palmada} \textsuperscript{(28)} o el \VUgras{volante} \textsuperscript{(29)}, que revolotea
de una \VUgras{raqueta} a la otra.

%%K t02-noto-1 noto
\VUnotes{143.2mm}{%
(*) Los juegos infantiles tienen a menudo nombres puramente nacionales.
Los hemos nombrado en ido lo mejor que hemos podido, ya inspirándonos en
la acción misma que los caracteriza, ya adoptando el nombre nacional de
uno u otro país cuando no resulta demasiado inexacto. Por ejemplo: el
\VUgras{juego del tonel} (alemán y francés) es el \VUgras{juego de la rana}
\textsuperscript{(20)} (inglés), en el que los jugadores tratan de echar sus
discos redondos por la boca de la rana de metal que está con la boca
abierta sobre la caja (el tonel) agujereada. El \VUgras{salto al hombro}
\textsuperscript{(30)} se distingue del \VUgras{salto de la rana} solo en esto:
para saltar por encima de la cabeza, los jugadores apoyan las manos en
los hombros del compañero, mientras que en el «\VUgras{salto de la rana}»
las apoyan solo en su espalda. --- O bien, algunos de los que juegan se
inclinan mientras los otros saltan sobre sus espaldas apoyando las manos
en los hombros; los primeros deben aguantar el golpe sin doblarse, los
segundos deben saltar sin perder el equilibrio (véase el cuadro mismo).%
}

%%K t02-26-1 p
26. --- En medio de la pradera hay dos árboles. Al pie de uno de ellos,
siete u ocho chicos han organizado una partida de \VUgras{salto al hombro}
\textsuperscript{(30)} (*). Este juego no se conoce en todos los países; pero
todo el mundo sabe qué es el \VUgras{salto de la rana}, al que los niños no
están jugando ahora.

%%K t02-tit-6 sub
\VUsaut{5.0mm}
\VUpk{10.2pt}{40}{Los Juegos al aire libre.}
\VUsaut{3.4mm}

%%K t02-27-1 p
27. --- La \VUgras{gallina ciega} \textsuperscript{(31)} es también muy divertida;
niñas y niños se ponen por turno la \VUgras{venda} en los ojos y atrapan a
los torpes que se dejan coger.

%%K t02-28-1 p
28. --- En medio de la pradera, he aquí un juego muy francés, aunque algo
brusco: es la \VUgras{osa} \textsuperscript{(32)}, mucho más ruidosa que el juego
tan tranquilo de la \VUgras{cometa} \textsuperscript{(33)} o que el juego inglés del
\VUgras{tenis} \textsuperscript{(34)}.

%%K t02-29-1 p
29. --- Este último deporte hace las delicias de los alumnos mayores.
Nuestros propios profesores lo practican de buena gana. Exige mucha
destreza y agilidad, y a mí me gusta muchísimo. Dejo el \VUgras{columpio}
\textsuperscript{(35)} para las niñas.

%%K t02-30-1 p
30. --- Tampoco me gusta otro juego inglés que puede llegar a ser
peligroso.

%%K t02-30-1b p
Me refiero al \VUgras{fútbol} \textsuperscript{(36)}, que hace la felicidad de mi
hermano mayor. Prefiero nuestro viejo \VUgras{juego de barras}
\textsuperscript{(38)}, igual de sano y mucho menos brutal.

%%K t02-tit-7 sub
\VUsaut{4.6mm}
\VUpk{10.2pt}{40}{Ciclismo y juegos de pelota.}
\VUsaut{3.0mm}

%%K t02-31-1 p
31. --- También doy con gusto una vuelta a la pradera en \VUgras{bicicleta}
\textsuperscript{(37)}.

%%K t02-32-1 p
32. --- El año que viene aprenderé a \VUgras{montar a caballo}
\textsuperscript{(39)}. ¡Qué hermoso es un caballo que anda o trota con garbo, y
que salta los obstáculos al galope!

%%K t02-33-1 p
33. --- En verano aprendemos a \VUgras{nadar} \textsuperscript{(40)}. Nos zambullimos
y nos bañamos con deleite en el agua clara y fresca del río. A veces, al
final, soltamos un \VUgras{globo} de papel \textsuperscript{(41)} que sube
majestuosamente por el aire en cuanto se llena de aire caliente.

%%K t02-34-1 p
34. --- Hay todavía muchos otros juegos, pero no acabaría nunca de
nombrarlos, y por este breve resumen se ve que los jueves en nuestra
hermosa pradera no son nada aburridos.

%%K t02-34-2 p
N. B. --- \textit{El profesor completará fácilmente este capítulo
describiendo con más detalle cada uno de los juegos importantes.}
\VUsaut{6.0mm}
\VUfilet{22mm}
